\documentclass[11pt]{article}
\usepackage{fullpage,graphicx,hyperref}

\newcommand{\fig}[2]{\begin{figure}[h]\begin{center}%
  \includegraphics[scale=#2]{#1.png}\end{center}%
  \end{figure}}
\newcommand{\ceil}[1]{\left\lceil #1\right\rceil}


\begin{document}

\begin{center}\Large\bf 
CS 473 (Spring 2025)\\
Homework 8 (due Apr 3 Thu 10am)
\end{center}

\medskip\noindent
{\bf Instructions:} As in previous homeworks.

\begin{description}

\bigskip
\item[Problem 7.1:]
Let $G=(V,E)$ be an unweighted graph with $n$ vertices and $m$ edges ($m\ge n$).
A \emph{loop} of length $\ell\ge 2$ refers to a walk $\langle v_1,v_2,\ldots,v_\ell,v_1\rangle$
such that $v_1,v_2,\ldots,v_\ell$ are all distinct and $v_1v_2,v_2v_3,\ldots,v_\ell v_1\in E$.  Note that a loop is the same as a (simple) cycle, except for the case of $\ell=2$ (usually $\langle v_1,v_2,v_1\rangle$ is not considered a cycle, but is a loop of length 2 if $v_1v_2\in E$).
We study the following problem: find a vertex-disjoint collection of loops and paths, maximizing their total length.

\begin{enumerate}
\item[(a)] (20 pts) Show that there always exists an optimal solution using only loops of length 2 or odd lengths, and no paths.  (So, we can redefine the problem without paths and the problem is the same.)  
\item[(b)] (10 pts) If $G$ is bipartite, show that the optimal total length is the same as $2k^*$, where $k^*$ is the maximum matching size.  (So, the problem is equivalent to maximum matching in the bipartite case.)
\item[(c)] (70 pts) Now consider the problem in the general case when $G$ is not bipartite.
Show that the problem can also be reduced to bipartite maximum matching and so
can be solved in $O(mn)$ time by the Hungarian method or $O(m\sqrt{n})$ time by
Hopcroft and Karp's algorithm.

(Hint: Construct a new bipartite graph $G'$ where we create two versions of each vertex and
two versions of each edge\ldots\ \ Carefully prove correctness.)
\end{enumerate}



\bigskip
\item[Problem 7.2:]
Given an undirected graph $G=(V,E)$ and an integer $d$, 
we want determine whether it is possible to direct the edges so that the resulting directed graph has maximum {\bf out-degree\/} at most $d$.
Describe how to solve this problem by reduction to maximum flow.
Prove correctness of your method.

Example: for the undirected graph below (left) and $d=2$, the answer is yes, and one solution is shown on the right (there are many
other solutions).

\fig{hw7fig2}{1}

(Hint: start with a bipartite graph where the vertices on the left side are \emph{edges} in $G$ and the vertices on the right side are the vertices in $G$;
add a source and a sink\ldots)



\bigskip
\item[Problem 7.3:]
We are given a set of $n$ axis-aligned rectangles $R=\{r_1,\ldots,r_n\}$, where each rectangle $r_i$ has a cost $a_i$.  The rectangles $r_i$ all have height~1 and have integer $x$- and $y$-coordinates, and no two rectangles intersect except along the boundaries, and their union is a rectangle $U(R)$.
We want to redraw the rectangles as $R'=\{r_1',\ldots,r_n'\}$, whose union
is a new rectangle $U(R')$, so that  the rectangles $r_i'$ still have
height~1 and have integer $x$- and $y$-coordinates, and 
the adjacency structure is unchanged, i.e., if a side of the
rectangle $r_i$ touches a side of rectangle $r_j$,
then the corresponding side of $r_i'$ touches the corresponding side of $r_j'$.

Describe a polynomial-time algorithm to compute a redrawing $S'$ that minimizes $\sum_{i=1}^n a_i\cdot\textrm{width}(r_i')$, where $\textrm{width}(r_i')$ denotes
the width of $r_i'$.

You may assume that there is a polynomial-time algorithm for the
following version of the \emph{minimum-cost flow} problem: given
a directed graph $G=(V,E)$ with source $s$ and sink $t$ and a value $d$, where
each edge $e\in E$ is given numbers $\ell(e),c(e),\textrm{cost}(e)$,
compute a flow $f$ of value $d$, minimizing $\sum_{e\in E} \textrm{cost}(e)\cdot f(e)$, such that for every $e\in E$, we have $\ell(e)\le f(e)\le c(e)$, and we have conservation of flow at every vertex in $V-\{s,t\}$.

(Hint: create a vertex for each rectangle\ldots) 

\fig{hw7fig3}{1}

\end{description}
\end{document}

